\begin{hongqiarticle}{
        title  = {深入持久地批判修正主义},
        author = {小靳庄大队党支部},
        style  = normal
    }

    听到伟大的领袖和导师毛主席逝世的消息，我们小靳庄的全体党员，贫下中
    农和社员群众，无不感到万分悲痛。大家思念着毛主席的恩情，止不住热泪滚
    滚，泣不成声。毛主席是我们的大救星，我们贫下中农永远怀念毛主席。

    我们一遍又一遍地学习中共中央、人大常委会、国务院，中央军委《告全党
    全军全国各族人民书》，学习华国锋同志在毛主席追掉大会上致的掉词，决心把悲
    痛化为巨大的力量。我们一定要继承毛主席的遗志，最紧密地团结在党中央的周
    围，遵照毛主席的嘱咐，“按既定方针办”，坚持以阶级斗争为纲，坚持党的基本路
    线，坚持无产阶级专政下的继续革命。当前，要继续深入开展批邓、反击右倾翻案
    风的斗争，坚持不解地批判修正主义，批判党内资产阶级，真正做到按毛主席的革
    命路线看问题，想事情，做工作，把毛主席开创的无产阶级革命事业进行到底。

    在批邓，反击右倾翻案风的斗争中，毛主席指出：“摘社会主义革命，不知道资
    产阶级在哪里，就在共产党内，党内走资本主义道路的当权派。走资派还在走。”
    党内走资派，为了颠覆无产阶级专政，复辟资本主义，总是打着红旗反红旗，推
    行修正主义路线，千方百计地反对毛主席的革命路线，反对以阶级斗争为纲。如
    果我们不批判修正主义，不同走资派斗，就分不清路线是非，就会在复杂的斗争
    中迷失方向，跟着修正主义路线跑。文化大革命前，我们小靳庄之所以成为全县
    有名的“老大难”单位，就是因为上面有刘少奇的修正主义路线，下面有阶级敌人
    和资本主义势力捣乱，党支部没有执行毛主席的革命路线，不抓阶级斗争，不批
    修正主义，为资本主义的泛滥开了绿灯。文化大革命以来，我们总结了过去的经
    验教训，狼批了刘少奇、林彪、邓小平的修正主义路线，带领广大群众大批资本主
    义，大干社会主义，迅速改变了小靳庄的落后面貌，成为一个农业学大寨的先进
    单位。在斗争中，我们不断增强了识别真假马克思主义的能力，提高了执行毛主席
    革命路线的自觉性。去年夏季前后，邓小平抛出“三项指示为纲”的修正主义纲
    领，大刮右倾翻案风，还恶毒攻击我们小靳庄。我们坚持用阶级斗争的观点观察
    问题，分析问题，看出了“三项指示为纲”反对以阶级斗争为纲，是同毛主席革命
    路线唱对台戏的。面对邓小平刮起的“十二级台风”，我们贫下中农“真理在胸旗
    在手，心红骨硬反潮流”，和他的修正主义路线对着干，坚持执行毛主席的革命
    路线。在社会主义时期，“走资派还在走”是一个长期存在的历史现象。我们一定
    要树立长期作战的思想，要继续批判邓小平反革命的修正主义路线，批判他授意
    炮制的《论总纲》、《汇报提纲》、《条例》等大毒草，把批判修正主义、批判资产阶
    级的斗争坚持下去。

    在批林批孔运动中，在同党内资产阶级的斗争中，我们越来越清楚地认识到，
    批判修正主义要同批判一切剥削阶级的意识形态结合起来。长期以来，剥削阶级
    的旧思想、旧文化、旧风俗、旧习惯，一直是束缚劳动人民的枷锁。在社会主义时期，
    剥削阶级虽然被推翻了，但是，他们的意识形态仍然顽固地存在着。正如列宁指
    出的，“旧社会灭亡的时候，它的死尸是不能装进棺材，埋入坟墓的。它在我们中
    间腐烂发臭并且毒害我们。”（全俄中央执行委员会，莫斯科工、农和红军代表苏维埃，工
    会联席会议》，《列宁全集》第27卷第407页）党内资产阶级为了推行修正主义路线，总
    是把反动的孔孟之道作为他们一个重要的思想支柱；社会上的阶级敌人也总是利
    用孔孟之道和一切腐朽没落的意识形态，向无产阶级进攻，作为他们搞复辟的一
    个重要手段。我们要执行毛主席的革命路线，用社会主义思想占领农村阵地，就
    必须深入开展意识形态领域里的社会主义革命，破旧立新，“在上层建筑其中包括
    各个文化领域中对资产阶级实行全面的专政”。文化大革命以来，我们狠抓了意识
    形态领域里的阶级斗争，带领广大群众不断扫荡剥削阶级的“四旧”，努力用社会主
    义的新思想、新文化占领农村阵地，并且在斗争中创办了“十件新事”。利用政治
    夜校、理论队伍、学唱革命样板戏、写革命诗歌，讲革命故事等多种形式，大造
    革命舆论，大批修正主义，大批资本主义，限制了资产阶级法权，巩固了农村社
    会主义阵地。邓小平之所以那样仇视小靳庄，拼命反对“十件新事”，就是因为我们
    这样做，夺了资产阶级的阵地，挖了修正主义的思想基础，使他的复辟阴谋难以
    得逞。这个事实也从反面告诉我们，抓好意识形态领域里的革命，对于反修防修，
    巩固无产阶级专政，防止资本主义复辟，建设社会主义，有着多么重大的意义。意
    识形态领域是阶级斗争的重要阵地。今后，我们一定要把这个方面的革命坚持下
    去，继续热情支持社会主义新生事物，发展“十件新事”，在毛主席的革命路线指引
    下，为造成使资产阶级既不能存在也不能再产生的条件而努力作战。

    毛主席指出：“阶级斗争是纲，其余都是目”。紧紧抓住了阶级斗争这个纲，
    坚持同修正主义斗，同党内走资派斗，就能取得革命和生产的主动权，什么困难
    都压不倒我们，而一定被我们所战胜。邓小平的“十二级台风”没有刮倒我们，七
    级地震的严重破坏，也阻挡不住我们继续革命的步伐，就是因为我们抓了阶级斗
    争，批了修正主义，批了资本主义。今年七月二十八日，我们小靳庄受到了强烈
    地震的破坏。震灾面前怎么办？我们反复学习了毛主席的一系列重要指示，紧紧
    抓住了阶级斗争这个纲。我们联系三年国民经济暂时困难时期两条路线斗争的历
    史经验和当前阶级斗争的表现，清醒地认识到，在存在阶级的社会中，人们同自
    然界的斗争总是和阶级斗争紧密联系在一起的，阶级敌人往往利用自然灾害造成
    的困难，进行揭乱和破坏。修正主义路线的头子，新老资产阶级都是这样干的。抗
    震救灾，既是同自然界的斗争，更是一场严重的阶级斗争。我们不能只见灾情，
    不见敌情，要严防阶级敌人趁机捣乱，严防修正主义，资本主义趁机泛滥。地震
    发生以来，我们坚持办好政治夜校，组织群众认真学习毛主席的重要指示和党中
    央的慰问电，联系实际，深入批邓，批判《论总纲》等三株大毒草，批判“白猫黑
    猫”论，牢固地树立人定胜天，靠社会主义战胜灾害的思想，鼓舞了广大群众夺
    取抗震救灾胜利的信心和决心。大家豪迈地说，“它来一次大地震，咱来一场大
    革命”。在批邓、反击右倾翻案风斗争的推动下，我们大队获得了又一个秋季大丰
    收。事实说明，不管是在顺利或在困难的情况下，都要坚持毛主席的革命路线，
    都要不停顿地开展对修正主义的批判。

    九月初，我们正准备给毛主席写报告，汇报我们大队如何在毛主席的革命路
    线指引下，深入开展批邓、反击右倾翻案风的斗争以及抗震救灾的工作。没想
    到，我们的报告还没有写完，毛主席他老人家就与世长辞了。毛主席啊，我们小靳
    庄人民永远怀念您，您是我们心中永远不落的红太阳。毛主席逝世了，战无不胜
    的毛泽东思想与世长存。我们决心永远坚持毛主席的革命路线，深入批邓、反击
    右倾翻案风，夺取抗震救灾的更大胜利，把同党内资产阶级的斗争进行到底，以
    实际行动把向毛主席的报告继续写下去，一直写到共产主义。
\end{hongqiarticle}